\section{ULA proofs}
\label{section:ULA}

\subsubsection{} The purpose of this section is to prove Lemma \ref{lem:zast_ULA}. In the case $P = B$, this follows from the work of J. Campbell \cite{justin_Kontsevich}. The main input to that work was the Kontsevich compactification, whose parabolic version $\Bunc_P^K \to \Bunc_P$ does not seem to lift along the map $\Bunt_P \to \Bunc_P$.

\subsection{Outline of the proof}

\subsubsection{} Let us write $\ZZ \mon \ZZt$ for $\ZZ_{N,P^-} \mon \ZZ_{\Nc, \Pt^-}$. We first reformulate the assertion of the lemma in terms that are smooth-local with respect to the data of
  \[j : \ZZ  \mon  \ZZt  \qquad  \qquad \FF = \left( \WW \shoximes j_{P,!}^- \right)_\ZZt.\]
So let us examine the purity morphism
    \[
    \left( \WW \stoximes j_{P,!}^- \right)_\ZZt 
    \to \left( \WW \shoximes j_{P,!}^- \right)_\ZZt \, [2 \delta_G]
    \]
in question. Note that the LHS is $!$-extended along $j$. Furthermore, the restriction of the purity map to $\ZZ$ is an isomorphism because $\WW \boxtimes j_{P,!}$ is a local system on $\Bun_N \times \Bun_{P^-}$. Therefore, our purity morphism identifies with the canonical morphism
\begin{equation}
\label{eqn:zast_ULA_rewrite}
    j_! \cdot j^! \cdot \FF \to \FF
\end{equation}
shifted by $[2 \delta_G]$. So Lemma \ref{lem:zast_ULA} is equivalent to the following assertion, which is smooth-local in the desired sense:

\begin{lemma}
\label{lemma:zast_!_extended}
The sheaf $\FF$ is $!$-extended from $\ZZ \mon \ZZt$.
\end{lemma}

\subsubsection{Proof of Lemma \ref{lemma:zast_!_extended}, modulo the rest of the section} Write $\UU \mon \ZZt$ for the largest open subset over which (\ref{eqn:zast_ULA_rewrite}) is an isomorphism. In Construction \ref{construction:zast_correspondence}, we will form a correspondence
\[
\begin{tikzcd}[
  column sep={4em,between origins},
  row sep={2em,between origins}
]
    & \HHt \arrow[dl, "\mu" above left] \arrow[dr, "\pi"] & \\
    \ZZt & & \ZZt
\end{tikzcd}
\]
whose action preserves $\UU$ in the sense that
    \[\mu^{-1} (\UU) = \pi^{-1} (\UU).\]
This last assertion is the content of Proposition \ref{proposition:zast_corr}. 

\subsubsection{Informal description of $\HHt$} This paragraph has no mathematical meaning. According to Proposition \ref{proposition:zast_loc_formula}, a point of $\ZZt$ over $D \in \Conf_{G,P}$ may be viewed as a degenerate $N_M$-bundle away from $D$ glued to some Zastava data on the formal neighborhood of $|D|$. Actually there is a constraint on the nature of this local Zastava data in terms of the value of $D$, but we will ignore it. In these terms, the action of $\HHt$ from right to left admits the following description:
\begin{itemize}
    \item Choose a divisor $D_\good$ whose support is disjoint from that of $D$, and impose the open condition that the $N_M$-bundle is nondegenerate at $D_\good$.
    \item Cut out the $N_M$-bundle at $D_\good$ and glue in some Zastava data coming from
        \[\Gr_N \cap \Gr_{P^-} \to \Gr_N \to \Gr_{N_M}.\]
    The result is a new point of $\ZZt$ lying over the divisor $D + D_\good$.
\end{itemize}

\subsubsection{} We will then prove Proposition \ref{proposition:SS_bound}, which gives an upper bound on the singular support of $j_{P,!}^-$. This will allow us to exhibit an open substack
    \[\ZZt_\Sing \mon \UU\]
with a fairly simple membership criterion. Finally, we prove Proposition \ref{proposition:Hecke_transitivity}, which says that that every point of $\ZZt$ can be related by the correspondence $\HHt$ to a point in $\ZZt_\Sing$.

\subsection{Construction of the correspondence}

\begin{construction} 
\label{construction:zast_correspondence}
There is a correspondence
\[
\begin{tikzcd}[
  column sep={4em,between origins},
  row sep={2em,between origins}
]
   & \HH \arrow[dd, hook] \arrow[dl] \arrow[dr] & \\
  \ZZ \arrow[dd, hook] &  & \ZZ \arrow[dd, hook] \\
   & \HHt \arrow[dl, "\mu" above] \arrow[dr, "\pi" above] & \\
  \ZZt & & \ZZt
\end{tikzcd}  
\]
with the feature that both squares are Cartesian.
\end{construction}

\begin{proposition}
\label{proposition:zast_corr}
The following hold:
\begin{itemize}
    \item[(i)] The map $\mu$ is etale. 
    \item[(ii)] The map $\pi$ is smooth-locally a product.
    \item[(ii)] The two pullbacks
      \[\mu^! \cdot \left(\WW \shoximes j_{P,!}^- \right)_\ZZt \qquad \pi^! \cdot \left(\WW \shoximes j_{P,!}^- \right)_\ZZt\]
    are smooth-locally isomorphic up to tensoring with the shift of a rank one local system.
\end{itemize}
\end{proposition}

\subsubsection{Notation} Our construction of $\HHt$ will be given in terms of the local description of $\ZZt$ recorded as Proposition \ref{proposition:zast_loc_formula}. We write
    \[
    \ZZ_\Conf = \ZZ_{N, P^-, \Conf}
    \qquad \ZZt_\Conf = \ZZ_{\Nc, \Pt^-, \Conf}.
    \]
The definition of $\ZZ_{N, P^-, \Conf}$ repeats the definition of $\ZZ_{\Nc, \Pt^-, \Conf}$ given in \S\ref{subsection:Zastava_spaces} but with $\Gr_N$ and $\Gr_{P^-}$ replacing $\Grc_N$ and $\Grt_{P^-}$. There is a natural map
    \[\ZZ_\Conf \to \Gr_{N,\Conf} \to \Gr_{N_M, \Conf}.\]

\subsubsection{} We take
  \[
  \HHt = \left\{ \Bunc_{N_M, \Conf \times \Conf}^{\vpol} \right\}_\disj \underset{[\Conf \times \Conf]_\disj}{\timest} \left\{ \ZZt_\Conf \times \ZZ_\Conf \right\}_\disj
  \]
Then the leftward map is
\begin{align*}
    \mu : \HHt \mon
    & \left\{ \Bunc_{N_M, \Conf \times \Conf}^{\vpol} \right\}_\disj \underset{[\Conf \times \Conf]_\disj}{\timest} \left\{ \ZZt_\Conf \times \ZZt_\Conf \right\}_\disj \\
    & \quad \simeq \ZZt \underset{\Conf}{\times} [\Conf \times \Conf]_\disj \\
    & \qquad \to \ZZt,
\end{align*}
where we are using the sum map
    \[
    m : [\Conf \times \Conf]_\disj \to \Conf
    \]
The rightward map is
\begin{align*}
   \pi : \HHt \to 
    & \left\{ \Bunc_{N_M, \Conf \times \Conf}^{\vpol} \right\}_\disj \underset{[\Conf \times \Conf]_\disj}{\timest} \left\{ \ZZt_\Conf \times \Gr_{N_M, \Conf} \right\}_\disj \\
    & \quad \simeq \left( \ZZt \underset{\Conf}{\times} [\Conf \times \Conf]_\disj \right)^{\text{good at $D_2$}} \\
    & \qquad \mon \ZZt \underset{\Conf}{\times} [\Conf \times \Conf]_\disj \\
    & \qquad \qquad \to \ZZt
\end{align*}
where we use the map
  \[\pr_1 : [\Conf \times \Conf]_\disj \to \Conf.\]
The base change of $\ZZ \mon \ZZt$ along either map identifies with
  \[\HH = \left\{ \Bunc_{N_M, \Conf \times \Conf}^{\vpol} \right\}_\disj \underset{[\Conf \times \Conf]_\disj}{\timest} \left\{ \ZZ_\Conf \times \ZZ_\Conf \right\}_\disj.\]

\subsubsection{} This completes Construction \ref{construction:zast_correspondence}. We now turn to the proof of Proposition \ref{proposition:zast_corr}.

\subsubsection{Proofs of (i) \& (ii)} These assertions are easily verified by examining the definitions. Indeed, $\mu$ is an open immersion followed by a base change of $m$, which is etale. This proves (i). To prove (ii), we note that twisted products are smooth-locally products and that $\pr_1$ is an open subset of a product.

\subsubsection{Proof of (iii)} This argument is very similar to the proof of Proposition \ref{proposition:Jac_Whit_preservation} appearing in \S\ref{subsection:Jacquet_Whittaker}. As before, we define $(\ZZt \times \Conf)^\good$ so that the lower square of the diagram
\[
\begin{tikzcd}[
  column sep={7.5em,between origins},
  row sep={4em,between origins}
]
  & \HHt \arrow[dl, dashed, "\text{$\widehat{\mu}$ or $\widehat{\pi}$}" above left] \arrow[d, "\text{$\mu$ or $\pi$}"] \\ 
  \left(\ZZt \times \Conf\right)^\good \arrow[r] \arrow[d] & \ZZt \arrow[d] \\
  \left( \Bunc_N \times \Conf \right)^\good \arrow[r] & \Bunc_N
\end{tikzcd}
\]
is Cartesian. As the dashed arrow indicates, both $\mu$ and $\pi$ have natural lifts fitting into a commutative diagram
\[
\begin{tikzcd}[
  column sep={8em,between origins},
  row sep={2em,between origins}
]
    & \HHt \arrow[dl, "\widehat{\mu}" above left] \arrow[dr, "\widehat{\pi}"] \arrow[dd, dashed] & \\
    \left(\ZZt \times \Conf \right)^\good \arrow[dd] & & \left(\ZZt \times \Conf \right)^\good \arrow[dd] \\
    & \HH_N \arrow[dl] \arrow[dr] \arrow[dd, "\psi"] & \\
    \left(\Bunc_N \times \Conf \right)^\good & & \left(\Bunc_N \times \Conf \right)^\good \\
    & \Ga &
\end{tikzcd}
\]
Therefore, the two $!$-pullbacks of $\WW$ to $\ZZt$ differ by the $!$-pullback of the Whittaker character along $\HHt \to \HH_N$. A similar argument shows that the two $!$-pullbacks of $j_{P,!}^-$ to $\ZZt$ are isomorphic.

\subsection{A singular support estimate}

\subsubsection{} The goal of this subsection is to obtain an upper bound on the singular support of $j_{P,!}^-$. In order to state the result, we will need to describe a stratification of $\Bunt_{P^-}$ into smooth stacks
  \[{}_{\vartheta' \le \vartheta}\Bunt_{P^-} \mon \Bunt_{P^-}\]
indexed by pairs of partitions.

\subsubsection{} Recall that $\Bunc_{P^-}$ has a smooth stratification by
  \[X^\theta \times \Bun_{P^-} \simeq {}_\theta\Bunc_{P^-} \mon \Bunc_{P^-}.\]
Under this identification. the map ${}_\theta\Bunc_{P^-} \to \Bun_G$ is given by projection on $\Bun_{P^-}$, followed by the usual map $\Bun_{P^-} \to \Bun_G$.

\subsubsection{} Taking the pullback along $\Bunt_{P^-} \to \Bunc_{P^-}$, we obtain a stratification by
    \[\Bun_{P^-} \underset{\Bun_M}{\times} \HH_M^{+, \theta} \simeq  {}_\theta\Bunt_{P^-} \mon \Bunt_{P^-}.\]
These strata are not smooth because $\HH_M^{+, \theta}$ is not smooth over $\Bun_M$, so let us examine $\HH_M^{+, \theta}$ more closely. Consider first the map
    \[\HH_M^{+, \theta} \to X^\theta.\]
We have a stratification $X^\vartheta \mon X^\theta$ indexed by partitions $\vartheta \vdash \theta$. The restriction $\HH_M^{+,\vartheta}$ of $\HH_M^{+, \theta}$ to each $X^\vartheta$ is a closed substack of $\HH_{M, X^\vartheta}$. Therefore, $\HH_M^{+, \vartheta}$ inherits a relative position stratification, indexed by $\vartheta' \le \vartheta$. These strata $\HH_M^{+, \vartheta' \le \vartheta}$ are smooth over $\Bun_M \times X^\vartheta$, and we define
    \[
    {}_{\vartheta' \le \vartheta} \Bunt_{P^-} = \Bun_{P^-} \underset{\Bun_M}{\times} \HH_M^{+, \vartheta \le \vartheta}.
    \]

\begin{proposition}
\label{proposition:SS_bound}
The singular support of $j_{P,!}^-$ is contained in the conormals of the strata
    \[
    {}_{\vartheta' \le \vartheta} \Bunt_{P^-}
    \mon \Bunt_{P^-}.
    \]
\end{proposition}

\subsubsection{} The proof of Proposition \ref{proposition:SS_bound} will use Zastava spaces as a local model for $\Bunt_{P^-}$, following \cite{BFGM}. Let us recall this. For the sake of notational consistency, we write
    \[
    \Mod_M^+ = \pt \underset{\Bun_M}{\times} \HH_M^+.
    \]
There are natural maps
    \[
    \ZZ_{U_P, \Pt^-} \to \Mod_M^+ \to \Conf_{G,P}.
    \]
and the first arrow has a canonical section. We write $(-)^\theta$ to denote base change to the component $\Conf^\theta \mon \Conf$. Now according to \cite[(1)]{BFGM_erratum}, the diagram of Zastava spaces on the left
\[
\begin{tikzcd}[
  column sep={1.5em},
  row sep={1em}
]
  \Mod_M^{+,\theta} \arrow[r, hook] \arrow[d, equal] & \ZZ_{U_P, \Pt^-}^\theta & \ZZ_{U_P, P^-}^\theta \arrow[l, hook', "j_\ZZ"] \\
  \Mod_M^{+,\theta} & & 
\end{tikzcd}
\qquad 
\begin{tikzcd}[
  column sep={1.5em},
  row sep={1em}
]
  {}_{\theta}\Bunt_{P^-} \arrow[r, hook] \arrow[d] & \Bunt_{P^-} & \Bun_{P^-} \arrow[l, hook'] \\
  \HH_M^{+,\theta} & & 
\end{tikzcd}
\]
is a smooth-local model for the diagram of stacks on the right. 
\subsubsection{Proof of Proposition \ref{proposition:SS_bound}} We make a preliminary construction and the perform a sequence of reductions.

\subsubsection{A group action} We claim that there is an action of $\LL_\Conf^+ M$ on $\ZZ_{U_P, \Pt^-}$ whose restriction to
    \[\Mod_M^{+, \vartheta} \mon \Gr_{M, X^\vartheta}\]
matches the usual action of $\LL_{X^\vartheta}^+ M$. This is because $\ZZ_{U_P, \Pt^-}$ admits a local description similar to that of $\ZZ_{N, \Pt^-}$. Namely, it is the base change of
    \[\Maps^\gen\left(\Disc, \pt \mon (U_P^- \bs G \git U_P) /M \right)\]
along the tautological section (\ref{eqn:tautological_section}). The left action of $M$ on $U_P^- \bs G \git U_P$ induces the action we seek. 

\subsubsection{First reduction} This model matches the strata
    \[\Mod_M^{+, \vartheta' \le \vartheta} \mon \Mod_M^{+, \theta}
    \qquad
    \HH_M^{+, \vartheta' \le \vartheta} \mon \HH_M^{+,\theta}.\]
Therefore, it suffices to show the singular support of $j_{\ZZ,!} \cdot k$ at every point of $\Mod_M^{+,\theta}$ is contained in the conormals to the strata $\Mod_M^{+, \vartheta' \le \vartheta}$.

\subsubsection{Second reduction} By the factorization of the Zastava diagram, we only need to consider points lying over the main diagonal $X \mon X^\theta$. These are the points belong to strata indexed by $\vartheta' \le (\theta)$ where $(\theta) \vdash \theta$ is the trivial partition.

\subsubsection{Third reduction} Since the Zastava diagram is etale local with respect to the curve, we may take $X = \AA^1$. Then the Zastava diagram admits a translation action by $\Ga$ commuting with the $\LL_\Conf^+ M$ action constructed above. The result follows because each stratum $\Mod_M^{+, \vartheta' \le (\theta)}$ is an orbit for the combined action of $\LL_\Conf^+ M \times \Ga$.

\subsection{End of the proof}

\subsubsection{} Let $\UU_\Sing \mon \Bunt_{P^-}$ be the open subset where the singular support of $j_{P,!}^-$ is transverse to the codifferential of $\Bunt_{P^-} \to \Bun_G$. The singular support bound implies:

\begin{proposition}
Let $x$ be a geometric point of ${}_\theta \Bunt_{P^-}$. If the image of $x$ under the projection ${}_\theta \Bunt_{P^-} \to \Bun_{P^-}$ belongs to the smooth locus of $\Bun_{P^-} \to \Bun_G$, then $x$ belongs to $\UU_\Sing$.
\end{proposition}

\subsubsection{Proof} We have the following Cartesian diagram:
\[
\begin{tikzcd}[
  column sep={1.5em},
  row sep={1.5em}
]
  \HH_M^{+, \vartheta' \le \vartheta} \arrow[r, hook] \arrow[dr, dashed] & \HH_M^{+, \theta} \arrow[d] & {}_\theta \Bunt_{P^-} \arrow[l] \arrow[d] & \\
  & \Bun_M & \Bun_{P^-} \arrow[l] \arrow[r] & \Bun_G
\end{tikzcd}
\]
Therefore it suffices to observe that diagonal dashed arrow is smooth.

\subsubsection{} Define
    \[\ZZt_\Sing = \UU_\Sing \underset{\Bunt_{P^-}}{\times} \ZZt.\]

\begin{proposition}
\label{proposition:Hecke_transitivity}
Every geometric point of $\ZZt$ can be taken to a point of $\ZZt_\Sing$ by the action of $\HHt$.
\end{proposition}

\subsubsection{Proof of Proposition \ref{proposition:Hecke_transitivity}} The proof comes in two parts.

\subsubsection{Part one} Let $z$ be a geometric point of $\ZZt$, and let $\FF_{P^-}$ be its image in $\Bun_{P^-}$. Then $\FF_{P^-}$ belongs to the smooth locus of $\Bun_{P^-} \to \Bun_G$ iff
\begin{equation}
\label{eqn:ULA_cohomology}  
    H^0\left(X, \FF_{P^-} \overset{P^-}{\times} \fg/\fp^- \right) \simeq 0
\end{equation}
To control this cohomology group, we choose an arbitrary reduction $\FF_{B^-}$ of $\FF_{P^-}$ to $B^-$ that is generically transverse to the image of $z$ in $\Bunc_N$. Then the vector bundle appearing in (\ref{eqn:ULA_cohomology}) acquires a complete flag whose graded pieces are the line bundles $\FF_T \overset{T}{\times} \alpha$ for $\alpha$ ranging over the roots appearing in $\fu_P$.

\subsubsection{Part two} From all of this data, we may extract the following defect divisors:
\begin{itemize}
    \item The image of $z$ in $\Conf_{G,P}$ is a divisor $D_1$, away from which $\FF_M$ has a degenerate $N_M$-reduction.
    \item The degeneracy of this $N_M$-reduction defines a $\Lambda_M$-valued divisor $D_2$ on $X \setminus |D_1|$.
    \item From the choice of $\FF_{B^-}$ we obtain a point of $\ZZ_{\overline{N},B^-}$ and hence a divisor $D_3$ in $\Conf_{G,B}$. 
\end{itemize}
Now consider a point $E$ of $\Conf_{G,P}$ whose support is disjoint from all of three of the above divisors, and let $E'$ be a lift of $E$ to $\Conf_{G,B}$. Perform a Hecke modification coming from
    \[
    \left(\ZZ_{N,B^-,\Conf}\right)_{E'} \to \left(\ZZ_{N,P^-,\Conf}\right)_E\]
to obtain a new point of $\ZZt$, say $z'$. The image of $z'$ in $\Bun_{P^-}$ comes with a reduction $\FF_{B^-}'$ to $B^-$, and we have
  \[\FF_T' \overset{T}{\times} \alpha \simeq \left(\FF_T \overset{T}{\times} \alpha \right)(\alpha(E'))\]
But we may choose $E$ so that the divisors $\alpha(E') = \alpha(E)$ for $\alpha$ appearing in $\fu_P$ are arbitrarily negative, thereby forcing the vanishing (\ref{eqn:ULA_cohomology}) for $\FF_{P^-}'$.